% Paper 3: CarbonStack — Unified Semiconductor AI Framework
% Target: DAC Workshop / ICCAD Demo

\documentclass[conference]{IEEEtran}
\usepackage{amsmath,graphicx,booktabs,hyperref}

\title{CarbonStack: One LLM Framework for\\Eight Semiconductor Design Tasks}

\author{
\IEEEauthorblockN{William Chen}
\IEEEauthorblockA{National Taiwan University}
}

\begin{document}
\maketitle

\begin{abstract}
We present CarbonStack, a unified framework that applies a single
local language model (Gemma4, 4.5B parameters) to eight semiconductor
design tasks: C++ optimization, HLS pragma prediction, code generation,
testbench generation, driver generation, simulation acceleration,
legacy modernization, and Python-to-C++ translation. All tasks run
fully offline with 100\% code safety via post-process merge.
Proof-of-concept validation shows 24/26 individual tasks successful
across all eight sub-projects, demonstrating that small MoE models
possess sufficient domain knowledge for diverse EDA automation.
\end{abstract}

\section{Introduction}

Semiconductor design involves a complex pipeline of tasks from
algorithm specification to physical implementation. Each stage
traditionally requires specialized tools and domain expertise.
Recent advances in large language models suggest potential for
automation, but existing approaches target individual tasks in
isolation.

We propose CarbonStack, an integrated framework where a single
4.5B-parameter model handles eight distinct design tasks. Unlike
cloud-based approaches (GPT-4, Claude), CarbonStack runs entirely
locally---critical for semiconductor IP protection.

\section{Framework Architecture}

CarbonStack shares a common core across all sub-projects:
\begin{itemize}
    \item \textbf{LLM Engine}: Gemma4 (4.5B active, MoE) via Ollama
    \item \textbf{Post-Process Merge}: Pragma/code extraction +
    safe merge back into original source
    \item \textbf{Verification}: Compiler-based validation
    (gcc, Vitis HLS, Verilator, iverilog)
\end{itemize}

\section{Eight Sub-Projects}

\subsection{CarbonCode: C++ Optimization}
Optimizes C++ code with HLS pragmas. PFS=0.834 on 30 ForgeHLS
kernels, validated by Vitis HLS with up to 63.3$\times$ speedup.

\subsection{Carbon-Verify: Testbench Generation}
Generates Cocotb/Verilog testbenches from module descriptions.
3/3 Cocotb testbenches generated, 3/3 Verilog compiled via iverilog.
Key finding: LLM testbenches verify code behavior, not spec compliance.

\subsection{Carbon-HLS: Code Generation}
Generates synthesizable HLS C from algorithm descriptions.
4/5 kernels generated, FFT validated on Vitis HLS (Fmax=137 MHz).

\subsection{Carbon-Firmware: Driver Generation}
Generates Linux kernel drivers from register specifications.
172-line platform driver with ioremap, request\_irq, file\_operations.

\subsection{Carbon-Sim: Simulation Acceleration}
Identifies computation-heavy sections in SystemVerilog BFMs and
suggests DPI-C replacements. 2/2 modules analyzed with correct
DPI-C import suggestions.

\subsection{Carbon-Legacy: Code Modernization}
Modernizes old C code (goto, macros, K\&R style) to C++17.
3/3 tasks: all goto removed, all macros replaced, const-correct.

\subsection{Carbon-Translate: Python to C++}
Translates Python algorithms to C++ with OpenMP parallelization.
4/4 translations successful, all include pragmas and const references.

\subsection{Carbon-Debug: Automatic Bug Detection}
Identifies and fixes common C/C++ bugs (buffer overflow,
use-after-free, integer overflow). 3/3 cases: average 2.3 bugs
found per case, all with fixes provided.

\section{Results Summary}

\begin{table}[t]
\centering
\caption{CarbonStack PoC results across 8 sub-projects.
All using Gemma4 (4.5B MoE) running locally.}
\label{tab:poc}
\begin{tabular}{llcc}
\toprule
\textbf{\#} & \textbf{Sub-Project} & \textbf{Tasks} & \textbf{Success} \\
\midrule
1 & CarbonCode (C++ opt)    & PIE + HLS  & 78\% pass@1 \\
2 & Carbon-Verify (testbench) & 3 modules  & 3/3 \\
3 & Carbon-HLS (code gen)    & 5 kernels  & 4/5 \\
4 & Carbon-Firmware (driver)  & 2 specs    & 1/1* \\
5 & Carbon-Sim (sim accel)    & 2 BFMs     & 2/2 \\
6 & Carbon-Legacy (modern.)   & 3 codes    & 3/3 \\
7 & Carbon-Translate (Py$\to$C++) & 4 scripts & 4/4 \\
8 & Carbon-Debug (bug fix)    & 3 bugs     & 3/3 \\
\midrule
  & \textbf{Total}           & \textbf{26} & \textbf{24/26} \\
\bottomrule
\end{tabular}
\vspace{1mm}
\footnotesize{*Firmware required max\_tokens=4096; default 2048 too short.}
\end{table}

\section{Discussion}

\subsection{One Model, Eight Tasks}
Gemma4's MoE architecture may explain its versatility: different
expert pathways activate for different task types (pragma optimization
vs. testbench generation vs. code translation).

\subsection{Telecom Relevance}
CarbonStack covers all core telecom FPGA kernels: FIR (5G filtering),
FFT (OFDM), LDPC (error correction), AES (encryption), Viterbi
(decoding). A 5G NR FIR filter was synthesized at 386 MHz on
Virtex UltraScale+.

\subsection{Limitations}
PoC-level validation only. Each sub-project needs deeper evaluation
(more modules, harder cases, real-world codebases) before industrial
deployment.

\section{Conclusion}

CarbonStack demonstrates that a single 4.5B local model can
meaningfully assist across the entire semiconductor design flow.
While no individual task reaches expert-level quality, the framework
provides a 80\% ``first draft'' across all tasks---saving engineers
hours of boilerplate work per module.

\bibliographystyle{IEEEtran}
\begin{thebibliography}{1}
\bibitem{chipnemo} M. Liu et al., ``ChipNeMo: Domain-Adapted LLMs for Chip Design,'' DAC, 2024.
\end{thebibliography}

\end{document}
